%% aLOTofImaginArches -- SoftwareX manuscript
%%
%% Build:  make          (in this directory)
%% Engine: pdflatex + bibtex, twice, as elsarticle expects.
%%
%% Every number in this paper is produced by the software itself: the figures
%% read .dat files written by paper/figures/makedata.js and makeexample.js, and
%% `make data` regenerates them. Nothing is transcribed by hand.

\documentclass[final,3p,times,twocolumn]{elsarticle}

\usepackage{amsmath,amssymb}
\usepackage{graphicx}
\usepackage{booktabs}

% elsarticle leaves essentially no space under a table caption, so booktabs
% draws \toprule across the caption's last line. Reproduced in a document with
% nothing but elsarticle and booktabs, so it is the class, not this paper.
\setlength{\belowcaptionskip}{6pt}
\usepackage{url}
\usepackage{pgfplots}
\pgfplotsset{compat=1.18}
\usepgfplotslibrary{fillbetween,groupplots}

% The palette of the application, so the figures of the paper and the figures
% on screen are recognisably the same drawing.
\definecolor{lotblue}{RGB}{0,114,189}
\definecolor{lotred}{RGB}{162,20,47}
\definecolor{lotbrick}{RGB}{196,156,110}
\definecolor{lotgrey}{RGB}{120,120,120}
\definecolor{lotgreen}{RGB}{60,140,60}

\pgfplotsset{
  lotaxis/.style={
    width=\linewidth, axis equal image, axis on top,
    tick align=inside, tick pos=both,
    grid=both, major grid style={lotgrey!20}, minor grid style={lotgrey!10},
    label style={font=\footnotesize}, tick label style={font=\scriptsize},
    legend style={font=\scriptsize, draw=lotgrey!50, fill=white,
                  fill opacity=0.85, text opacity=1},
  },
}

\journal{SoftwareX}

\begin{document}

\begin{frontmatter}

\title{aLOTofImaginArches: interactive graphical statics of masonry arches
       in the browser}

\author{G.~Castellazzi}
\address{Department of Civil, Chemical, Environmental and Materials
  Engineering (DICAM), University of Bologna, Viale Risorgimento 2,
  40136 Bologna, Italy \\
  \small{\tt giovanni.castellazzi@unibo.it}}

\begin{abstract}
The equilibrium of a masonry arch is not a single state but a
three-parameter family of them, and that multiplicity is the hardest idea to
convey when the subject is taught. \emph{aLOTofImaginArches} makes it
manipulable. A student loads a photograph of a real arch, traces its intrados
and extrados, subdivides the ring into voussoirs, sets a scale, applies loads,
and then moves the three degrees of freedom of the line of thrust while the
software reports, joint by joint, whether the line stays within the masonry:
Heyman's safe theorem, made operable. The software is a rewrite of an
established MATLAB application as dependency-free ES modules that run directly
from a static host, so that no licence and no installation stand between a
student and the construction. The graphical idiom of the original is preserved
deliberately. Twenty-seven legacy examples are carried over and fifteen of
them, which store a complete solution, are reproduced to machine precision.
\end{abstract}

\begin{keyword}
masonry arches \sep line of thrust \sep graphical statics \sep limit analysis
\sep engineering education \sep web application
\end{keyword}

\end{frontmatter}

%% ---------------------------------------------------------------- metadata --

\begin{table*}[t]
\footnotesize
\centering
\caption{Code metadata}
\label{tab:metadata}

\begin{tabular}{@{}p{0.04\textwidth}p{0.36\textwidth}p{0.52\textwidth}@{}}
\toprule
\textbf{Nr.} & \textbf{Code metadata description} & \textbf{Please fill in this column} \\
\midrule
C1 & Current code version & v1.0.0 \\
C2 & Permanent link to code/repository used for this code version &
     \url{https://github.com/gcastellazzi/aLOTofImaginArches} \\
C3 & Permanent link to Reproducible Capsule &
     The application itself is the capsule: it is served as static files and
     runs, unmodified, at
     \url{https://gcastellazzi.github.io/aLOTofImaginArches/app/} \\
C4 & Legal Code License & MIT License \\
C5 & Code versioning system used & git \\
C6 & Software code languages, tools, and services used &
     JavaScript (ES2022 modules), HTML5 Canvas, CSS; Python~3 for the
     legacy-format converter; GitHub Pages \\
C7 & Compilation requirements, operating environments, dependencies &
     None. Any modern browser runs the application as served; there is no
     build step and no third-party runtime dependency. Node.js~$\ge$~18 is
     needed only to run the test suite; Python~3 with NumPy and SciPy only to
     re-run the converter. \\
C8 & Link to developer documentation/manual &
     \url{https://gcastellazzi.github.io/aLOTofImaginArches/} and the
     \texttt{README.md} files under \texttt{docs/app/} \\
C9 & Support email for questions & giovanni.castellazzi@unibo.it \\
\bottomrule
\end{tabular}
\end{table*}

%% --------------------------------------------------- 1. Motivation ---------

\section{Motivation and significance}

A masonry arch stands not because it is in \emph{the} state of equilibrium but
because it is in \emph{one} of infinitely many. Under Heyman's assumptions ---
no tensile strength, unlimited compressive strength, no sliding ---
the safe theorem states that if any line of thrust in equilibrium with the
applied loads can be found lying wholly within the masonry, the structure is
safe \cite{Kooharian1952,Heyman1966,Heyman1982}. For a two-dimensional arch
that line has three degrees of freedom, the $\infty^3$ of the classical
literature \cite{Heyman1995,Huerta2001}.

Students meet this in a lecture and, in the author's experience of teaching it,
do not believe it. The obstacle is not the algebra of the funicular polygon,
which is elementary, but the multiplicity: a single worked example produces a
single drawing, and a single drawing looks like a single answer. What convinces
is moving the solution and watching it remain a solution --- and then watching
it stop being an \emph{admissible} one as it leaves the ring.

Interactive equilibrium tools have a distinguished record in exactly this role
\cite{ODwyer1999,Block2006,Block2007}, and this software follows that line for
the two-dimensional arch, with an emphasis on the historical construction
rather than on generality. Its predecessor, a MATLAB App used for several years
in the master's courses \emph{Mechanics of Historical Masonry Structures} and
\emph{Historical Masonry and Wooden Structures} at Bologna, was limited by two
things: it required a MATLAB licence, which excludes students working on their
own machines and anyone outside a subscribing institution, and its source,
packaged in a binary \texttt{.mlapp} container, was not something a student
could read.

This paper describes a rewrite as a browser application: plain ES modules with
no dependencies and no build step, served as static files, so that the source
a student reads is exactly the source that runs. Three commitments shaped it.

\begin{enumerate}
\item \textbf{The graphics of the original are preserved deliberately.} The
  boxed axes with inward ticks, the pale grid, the translucent force bands and
  the constant-size arrowheads of the MATLAB figure are reproduced on the HTML
  canvas. Continuity of appearance is not nostalgia: the course notes, the
  slides and the existing tutorials all show that idiom, and a student
  switching between them should not have to re-learn what a picture means.
\item \textbf{The mechanics is separated from the drawing and is tested.} The
  construction lives in seven small modules with no reference to the DOM, so
  it can be, and is, exercised by a test suite from the command line.
\item \textbf{The legacy examples are carried over, and checked.} Every stored
  example of the MATLAB version was converted, and the conversion was verified
  field by field rather than assumed.
\end{enumerate}

%% ------------------------------------------------- 2. Software description --

\section{Software description}

\subsection{Software architecture}

The application is a static web page. There is no server, no bundler and no
package to install: the browser loads \texttt{index.html}, which imports ES
modules directly.

The mechanics lives in nine \texttt{core} modules -- geometry, blocks,
statics, tracing, units, the saved-model reader, the save format, the
mechanism and the dome -- with three more for the drawing and one for the
interface: about 4\,500 lines in all.

The separation is strict in one direction: \texttt{core} knows nothing of the
canvas, and \texttt{render} knows nothing of arches. It is what allows the
mechanics to be checked without a browser. The suite comprises 123 tests run by
the Node test runner (\texttt{node --test}), covering the geometry, the
funicular construction, the reproduction of stored examples, tracing, scale and
units, admissibility, hinge formation and the collapse kinematics, the dome,
and the save format.

A separate Python tool, \texttt{tools/mat2json.py}, converts the MATLAB
\texttt{.mat} examples to JSON and extracts the background images. It carries
a \texttt{--check} mode that re-reads every converted file and compares it
field by field with the original, so the conversion is verified rather than
trusted.

\subsection{Software functionalities}

\paragraph{Stored examples} Twenty-eight example files are shipped. One is an
empty template; of the remaining twenty-seven, fifteen carry a complete stored
solution and are recomputed from their geometry alone, six were saved before a
solution was computed, and six are internally inconsistent --- the stored
solution does not correspond to the stored geometry. The software detects the
last group and says so, rather than drawing a solution it cannot justify
(Section~\ref{sec:legacy}).

\paragraph{Tracing} A student loads any image and clicks along the intrados and
the extrados. The two curves are resampled by arc length and joined into $n$
voussoirs with $n+1$ radial joints. The trace is checked before it is used: too
few points, coincident curves, and --- the case that matters --- curves that
cross are all reported. A crossing is detected by the \emph{sign} of the signed
area of the resulting blocks, not by their size, because an inverted block can
be as large as a sound one.

\paragraph{Scale, units and loads} Until a scale is set, every quantity is in
pixels. The student picks two points and states the real distance between them;
lengths then scale by $k$, areas by $k^2$, and weights by $k^2$ or $k^3$
according to whether the out-of-plane thickness was given physically or in
pixels --- a distinction worth making explicit, since getting it wrong is a
silent factor of the scale. Three unit systems are offered, matching the
original. Span and rise are reported after scaling, being the quickest check
that the reference distance was entered correctly. Point loads are placed by
clicking; a load and a voussoir weight are stations on the load line alike, so
loads are merged into the same ordered sequence as zero-area blocks.

\paragraph{The three degrees of freedom} The pole of the force polygon has two
coordinates -- its abscissa is the horizontal thrust, its ordinate divides the
total weight between the two vertical reactions -- and the line has one free
end, sliding along its springing joint from the intrados to the extrados. Each
is a slider. The far end is not a fourth parameter: it follows, because the
last ray is carried on until it meets the other joint, and where it lands is
reported.

\paragraph{Admissibility} For each joint the software computes where the line
crosses it, as a fraction $s$ of the joint. The verdict names the number of
joints violated, the worst one and on which face, or --- when the line fits ---
how much room is left. Choosing the crossing correctly is less trivial than it
appears: a joint near a springing is almost horizontal, and its infinite line
meets segments of the thrust line far away that have nothing to do with it, so
the implementation prefers the crossing that lies \emph{within} the joint.

\paragraph{Mechanism} The line cannot leave the ring: where it would, it stays
hooked at the intrados or the extrados between two voussoirs, and that point is
a hinge. Counting the two springings as hinges throughout, $h$ hinges carry
$h-1$ rigid macro-blocks and the arch has $3(h-1)-2h = h-3$ degrees of freedom:
two hinges is the once-hyperstatic arch whose equilibrium state is not
determined, three is the three-pin arch, four is a mechanism. In this mode the
thrust slider alone commands the line, the other two parameters being chosen to
hold it as far from both faces as it will go, so that hinges appear only when
the thrust forces them. The collapse kinematics is drawn
(Section~\ref{sec:mechanism}).

\paragraph{Poleni's dome} A barrel vault cut into voussoirs gives blocks of
constant width. A dome cut into lunes does not: each lune is bounded by two
meridian planes, so its width is proportional to the distance from the axis,
broad at the major parallel and closing to nothing at the crown. With the
option set, the weight of each voussoir is taken from Pappus' theorem,
$V = A\,\theta\,\bar r$ with $\bar r$ the distance of the block's centroid from
the axis --- exact for a plane region turned about an axis in its plane. The
axis defaults to the mid-point of the springings and can be picked on the
drawing. A second view of the right-hand pane shows the voussoirs as solids,
prismatic or revolved.

\paragraph{Hooke's cable} The same polygon reflected about the chord joining
the two springings, so that the cable hangs from the points the arch springs
from whatever their relative height, and \emph{ut pendet continuum flexile,
sic stabit contiguum rigidum inversum} \cite{Hooke1676} can be shown rather
than quoted. The rays of the force polygon are drawn dashed and lettered after
Bow, with the same letter on the parallel segment of the thrust line, so that a
force can be carried from one drawing to the other by eye.

\paragraph{Saving} The whole session --- traced curves, blocks, joints,
weights, loads, scale, unit system, dome settings and the three parameters ---
is written to a single JSON file and read back. Everything derivable is
deliberately omitted and recomputed on opening, so a file stays readable when
the mechanics is corrected; a file is validated before it is trusted and
refused with a diagnostic rather than half-loaded.

%% ------------------------------------------------ 3. Illustrative examples --

\section{Illustrative examples}

\subsection{Reproducing Poleni}
\label{sec:legacy}

Figure~\ref{fig:poleni} shows the software working over Poleni's own plate of
1748 for the dome of St~Peter's \cite{Poleni1748}, the drawing in which the
hanging-chain analogy was first put to a real structure. The 56 voussoirs and
their weights come from the stored example; the force polygon and the line of
thrust are recomputed from that geometry by the JavaScript implementation. The
largest discrepancy between the recomputed line and the one MATLAB saved is
$8.0\times10^{-16}$ relative --- machine precision. Across all fifteen
examples that carry a complete solution, the worst relative discrepancy is
$8.5\times10^{-16}$ for the force polygon and $3.9\times10^{-15}$ for the line
of thrust.

That agreement is worth stating precisely because the remaining examples do
\emph{not} agree, and for identifiable reasons: in some the stored force
polygon was drawn with weights that no longer match the stored blocks; in
others, applied loads were merged into the sequence but the loads themselves
were never written to the file. This is a property of the original save format,
not of either implementation, and it was found only because the port was
checked against the original rather than eyeballed. The software now reports
these files as not recomputable and displays the stored solution with a
warning.

\subsubsection*{The dome Poleni was actually analysing}

Reproducing the stored file exercises the port; it does not exercise the
mechanics Poleni was after. The plate is a section of a \emph{dome}, and a
lune of a dome is not a slice of a barrel: its width grows with the distance
from the axis. Setting the Poleni option on the same 56 voussoirs, with the
axis where the software places it by default, moves \textbf{14.7 per cent of
the total weight} from the upper courses to the lower ones. The outermost block
rises from 4.15 to 5.56 per cent of the total and the innermost falls from 1.29
to 0.34. The line of thrust follows, and the two idealisations of the same
drawing part company --- which is the whole of Poleni's argument, and something
the constant-thickness reading cannot show.

\begin{figure*}[t]
\centering
% Figure: the Poleni example, recomputed.
%
% The plate is drawn as a background image and everything on top of it is
% vector graphics read from .dat files written by makeexample.js. The image on
% disk was reduced when the examples were converted, but the model coordinates
% refer to the original plate, so the axis limits are the ORIGINAL pixel size
% and the reduced image is stretched onto them -- the aspect ratio is the same,
% so nothing is distorted.
%
% poleni_plate.jpg is the converted plate flipped about its horizontal axis:
% row 0 of a photograph is its top, which with an upward axis is the bottom of
% the box. The application performs the same reflection on the canvas.
\def\exName{Poleni\_Example\_01}
\def\exBlocks{56}
\def\exImage{Poleni_Example_01.jpg}
\def\exImgW{1767}
\def\exImgH{2161}
\def\exThrust{1.975e+8}
\def\exPoleX{197500000.000}
\def\exPoleY{-873452294.557}
\def\exRelErr{8.0e-16}


\begin{tikzpicture}
\begin{groupplot}[
  group style={group size=2 by 1, horizontal sep=1.1cm},
  lotaxis, height=7.2cm,
]

\nextgroupplot[
  width=0.62\textwidth, axis equal image,
  xlabel={$x$ [px]}, ylabel={$y$ [px]},
  xmin=0, xmax=\exImgW, ymin=900, ymax=\exImgH,
  grid=none,
  title={\footnotesize Poleni's plate, 1748, with the recomputed model},
]
\addplot graphics [xmin=0, xmax=\exImgW, ymin=0, ymax=\exImgH]
  {figures/poleni_plate.jpg};

% The 56 voussoirs, as closed paths separated by nan rows.
\addplot [lotgrey, line width=0.25pt, fill=lotbrick, fill opacity=0.20]
  table [x=x, y=y] {figures/data/Poleni_Example_01_blocks.dat};

% Hooke's cable: the same polygon, inverted.
\addplot [lotgreen, line width=0.9pt, dash pattern=on 3pt off 2pt]
  table [x=x, y=y] {figures/data/Poleni_Example_01_cable.dat};

% The line of thrust.
\addplot [lotblue, line width=1.3pt]
  table [x=x, y=y] {figures/data/Poleni_Example_01_lot.dat};

\legend{}

\nextgroupplot[
  width=0.30\textwidth, axis equal=false,
  xlabel={$x$}, ylabel={$y$},
  scaled ticks=true,
  title={\footnotesize Force polygon},
  grid=none,
]
% The rays, from the pole to every division of the load line.
\addplot [lotgrey, line width=0.2pt]
  table [x=x, y=y] {figures/data/Poleni_Example_01_rays.dat};
% The load line itself.
\addplot [lotred, line width=1.1pt]
  table [x=x, y=y] {figures/data/Poleni_Example_01_loadline.dat};
% The pole.
\addplot [only marks, mark=*, mark size=1.3pt, black]
  coordinates {(\exPoleX, \exPoleY)};
\node[anchor=west, font=\scriptsize] at (axis cs:\exPoleX,\exPoleY) {\,$O$};

\end{groupplot}
\end{tikzpicture}

\caption{Poleni's plate of 1748 with the model recomputed by the software:
  56 voussoirs (grey), the line of thrust (blue) and the inverted Hooke cable
  (green), with the force polygon on the right. The line is recomputed from
  the geometry alone and agrees with the solution stored by the MATLAB
  application to $8.0\times10^{-16}$ relative.}
\label{fig:poleni}
\end{figure*}

\subsection{The minimum thickness of a semicircular arch}

The second example is the classical one, and it exposed a defect worth
recording. For a semicircular ring of uniform thickness the least $t/r_i$ at
which an admissible line exists is about $0.108$ \cite{Heyman1982,Huerta2006}.

An early version of this software pinned both ends of the line at the
mid-points of the end joints, leaving only the thrust free. Measured by
bisection on a ring of 16 voussoirs, that one-parameter family first admits a
line at $t/r_i = 0.198$ --- nearly twice Heyman's value. The tool was
therefore condemning arches that stand.

Two things were wrong. The ends were pinned, discarding two degrees of freedom;
and, less obviously, the funicular construction inherited from the original
stops on the \emph{vertical} through the far springing rather than reaching the
far joint at all, so asking where the line crossed that joint returned the
intersection of the joint's infinite line with an unrelated segment. Freeing
the start along its joint and carrying the last ray on to the far joint
restores the full family. Table~\ref{tab:minthick} gives the measured result.

\begin{table}[t]
\footnotesize
\centering
\caption{Least $t/r_i$ admitting a line of thrust on a semicircular ring, by
  bisection, against the number of voussoirs $n$. Heyman's value for a
  continuous ring is $0.108$.}
\label{tab:minthick}

% Generated by paper/figures/makedata.js -- do not edit.
\begin{tabular}{@{}rcc@{}}
\toprule
$n$ & ends pinned & ends free \\
\midrule
8 & 0.1789 & 0.1116 \\
12 & 0.1894 & 0.1118 \\
16 & 0.1984 & 0.1155 \\
24 & 0.2070 & 0.1181 \\
32 & 0.2061 & 0.1179 \\
48 & 0.2085 & 0.1199 \\
64 & 0.2093 & 0.1202 \\
\bottomrule
\end{tabular}

\end{table}

With the ends free the figure settles at about $0.12$ and approaches Heyman's
$0.108$ from above; the residual gap is discretisation, a funicular polygon of
$n$ segments being able to cut corners that a continuous curve cannot, and it
narrows as $n$ falls rather than rises for exactly that reason. More telling
than the number is that the search, given no guidance, recovers the textbook
limit state on its own: the critical line leaves one springing at the
\emph{extrados}, arrives at the other at the \emph{extrados}, and requires a
horizontal thrust of about one fifth of the total weight.

Figure~\ref{fig:band} gives the admissible range of thrust against thickness
for both families: the safe theorem, drawn. The free band opens at
$t/r_i \approx 0.12$ and has already widened to $0.190 \le H/W \le 0.215$ at
$0.14$, a thickness at which the pinned family still admits nothing at all; the
pinned band does not open until $0.20$.

\begin{figure}[t]
\centering
% Figure: the admissible thrust band against thickness, pinned and free.
%
% Data from makedata.js: for each t/ri, the least and greatest H/W for which
% SOME admissible line exists in each family.
\begin{tikzpicture}
\begin{axis}[
  lotaxis, width=0.92\linewidth, height=5.4cm, axis equal image=false,
  xlabel={thickness ratio $t/r_i$},
  ylabel={horizontal thrust $H/W$},
  xmin=0.08, xmax=0.6, ymin=0.04, ymax=0.45,
  legend pos=north west, legend cell align=left,
]
\addplot [name path=fA, draw=none, forget plot]
  table [x=tri, y=Hmin] {figures/data/band_free.dat};
\addplot [name path=fB, draw=none, forget plot]
  table [x=tri, y=Hmax] {figures/data/band_free.dat};
\addplot [lotblue!18] fill between [of=fA and fB];
\addlegendentry{ends free}

\addplot [name path=pA, draw=none, forget plot]
  table [x=tri, y=Hmin] {figures/data/band_pinned.dat};
\addplot [name path=pB, draw=none, forget plot]
  table [x=tri, y=Hmax] {figures/data/band_pinned.dat};
\addplot [lotred!25] fill between [of=pA and pB];
\addlegendentry{ends pinned at the joint mid-points}

\addplot [lotblue, line width=0.9pt, forget plot]
  table [x=tri, y=Hmin] {figures/data/band_free.dat};
\addplot [lotblue, line width=0.9pt, forget plot]
  table [x=tri, y=Hmax] {figures/data/band_free.dat};
\addplot [lotred, line width=0.9pt, forget plot]
  table [x=tri, y=Hmin] {figures/data/band_pinned.dat};
\addplot [lotred, line width=0.9pt, forget plot]
  table [x=tri, y=Hmax] {figures/data/band_pinned.dat};
\end{axis}
\end{tikzpicture}

\caption{Range of horizontal thrust, as a fraction of the total weight, for
  which an admissible line exists, against the thickness ratio. Freeing the
  ends both lowers the thickness at which the band opens and widens it
  everywhere.}
\label{fig:band}
\end{figure}

\subsection{Hinges, and what a degree of freedom does not tell you}
\label{sec:mechanism}

Driven from the thrust, the same ring produces the classical collapse patterns
without being told what to look for. At $t/r_i = 0.18$ the arch stands for
$0.171 \le H/W \le 0.238$. At the lower limit five hinges have formed --- the
two springings, the intrados at the haunches and the extrados at the crown ---
giving four macro-blocks and two degrees of freedom; at the upper limit, four
hinges and one. In between, the line runs clear of both faces, nothing is
located, and the arch is once hyperstatic: its equilibrium state is not
determined, which is the $\infty^3$ restated in the language of constraints.

Figure~\ref{fig:mechanism} shows the minimum-thrust mechanism displaced. Two
things about that drawing were harder than they look.

The first is that an instantaneous centre is instantaneous. Turning each
macro-block about its centre through a \emph{finite} angle keeps the blocks
rigid but opens the hinges at second order --- on a chain of span 8, an
amplitude of $0.2$ radians parts them by $0.14$, which is plainly visible. The
motion is therefore integrated: at every step the velocity field is re-solved
for the current hinge positions and each block advanced a little about its
current centre, so every block undergoes a composition of rigid transforms and
stays exactly rigid while the residual opening falls as
$\textrm{amplitude}^2/\textrm{steps}$, to $1.2\times10^{-3}$ at the default.

The second is that \textbf{masonry cannot pass through itself}, and the sign of
a null-space vector falls out of the elimination rather than out of the
mechanics. A hinge sits at one face of its joint and the joint must open at the
other, so the sense of the motion is taken from the joints, not from the
arithmetic. Some hinge patterns admit \emph{no} good sense: on a symmetric ring
at maximum thrust both haunch hinges fall on the intrados, which leaves the
crown block turning about a point on the axis instead of dropping, and one
haunch opens exactly as the other shuts. The software reports that the pattern
would need the masonry to interpenetrate rather than animating it. \emph{A
positive degree of freedom is not by itself a collapse mode}: the joints have
to be able to open, and the count alone does not say whether they can.

\begin{figure}[t]
\centering
% Figure: the five-hinge minimum-thrust collapse mechanism.
%
% Data from makedata.js: the ring at rest, the same blocks displaced by the
% integrated motion, and the hinges. Every joint opens -- the amounts are in
% the log of the generator and quoted in the caption.
\def\mechHinges{5}
\def\mechBodies{4}
\def\mechDof{2}
\def\mechThrust{0.171}
\def\mechPattern{SieiS}


\begin{tikzpicture}
\begin{axis}[
  lotaxis, width=0.92\linewidth, height=5.0cm,
  xlabel={$x/r_i$}, ylabel={$y/r_i$},
  xmin=-1.32, xmax=1.32, ymin=-0.10, ymax=1.34,
  grid=none, ytick={0,0.5,1},
  legend style={at={(0.5,-0.30)}, anchor=north, draw=lotgrey!50,
                font=\scriptsize},
  legend cell align=left,
]
\addplot [lotgrey!70, line width=0.35pt]
  table [x=x, y=y] {figures/data/mech_rest.dat};
\addlegendentry{the arch at rest}

\addplot [lotred, line width=0.8pt]
  table [x=x, y=y] {figures/data/mech_moved.dat};
\addlegendentry{the mechanism, displaced}

\addplot [only marks, mark=*, mark size=2.1pt, lotblue]
  table [x=x, y=y] {figures/data/mech_hinges.dat};
\addlegendentry{hinges}
\end{axis}
\end{tikzpicture}

\caption{The minimum-thrust collapse mechanism of a semicircular ring,
  $t/r_i = 0.18$, $n = 16$, at $H/W = \mechThrust$: \mechHinges{} hinges,
  \mechBodies{} rigid macro-blocks, \mechDof{} degrees of freedom. The
  displacement is integrated rather than applied as a single rotation, so the
  blocks stay rigid and every joint opens instead of interpenetrating.}
\label{fig:mechanism}
\end{figure}

%% --------------------------------------------------------------- 4. Impact --

\section{Impact}

The immediate impact is on teaching. The software is used in the master's
course \emph{Mechanics of Historical Masonry Structures} at the University of
Bologna, alongside a set of lecture slides, handouts and exercises derived from
\cite{CastellazziTaliercio2025}. Removing the licence requirement changes what
can be asked of students: an exercise can now be set as \emph{photograph an
arch in this city, trace it, and find an admissible line}, and every student
can do it on their own machine, at home, without an installation. The saved
session is a single small JSON file, which makes such an exercise collectable
and markable.

A second effect is on reproducibility. Because the port was verified against
the original rather than reimplemented from its description, it produced a
usable audit of a decade of stored examples and identified six that cannot be
reproduced from what they contain --- a researcher relying on them would have
had no warning. That a rewrite is an opportunity to audit the data of the
original, and that the audit should be automated and kept, applies well beyond
this software. Third, the code is small enough to be read: a 300-line module in
which \texttt{forcePolygon} and \texttt{funicular} sit side by side with the
classical terms as their variable names is itself a teaching artefact.

The limitations follow from Heyman's assumptions rather than from the
implementation: the analysis is two-dimensional, blocks are rigid, sliding is
not checked, and the finite subdivision makes the computed line a polygon
inscribed in the true funicular curve. None are hidden from the user; the
discretisation, in particular, is a slider.

%% ---------------------------------------------------------- 5. Conclusions --

\section{Conclusions}

\emph{aLOTofImaginArches} makes the three-parameter family of equilibrium
states of a masonry arch something a student can take hold of and move. The
rewrite from MATLAB to dependency-free browser modules removes the licence and
the installation, preserves the graphical language of the original, and --- by
separating the mechanics from the drawing --- makes the construction testable
from the command line. The legacy examples are reproduced to machine precision
where they can be reproduced at all, and the software says so when they cannot.

Freeing the two end positions of the line, which the original construction had
fixed, brought the computed minimum thickness of a semicircular ring from
$t/r_i = 0.198$ to $0.115$ against Heyman's $0.108$, and let the software
recover the classical limit state without being told what to look for. The same
holds of the collapse patterns, which the hinge analysis finds unaided, and of
Poleni's own question: treating his plate as the dome it is, rather than as a
barrel of the same profile, redistributes a seventh of the weight. Work in
progress extends the treatment to arches on spreading supports and to the
sliding condition that Heyman's third assumption sets aside.

\section*{Acknowledgements}

The MATLAB predecessor was developed for, and improved by, successive cohorts
of students of the courses \emph{Mechanics of Historical Masonry Structures}
and \emph{Historical Masonry and Wooden Structures} at the University of
Bologna.

\bibliographystyle{elsarticle-num}
\bibliography{refs}

\end{document}
